% Options for packages loaded elsewhere
\PassOptionsToPackage{unicode}{hyperref}
\PassOptionsToPackage{hyphens}{url}
\documentclass[
  ignorenonframetext,
]{beamer}
\newif\ifbibliography
\usepackage{pgfpages}
\setbeamertemplate{caption}[numbered]
\setbeamertemplate{caption label separator}{: }
\setbeamercolor{caption name}{fg=normal text.fg}
\beamertemplatenavigationsymbolsempty
% remove section numbering
\setbeamertemplate{part page}{
  \centering
  \begin{beamercolorbox}[sep=16pt,center]{part title}
    \usebeamerfont{part title}\insertpart\par
  \end{beamercolorbox}
}
\setbeamertemplate{section page}{
  \centering
  \begin{beamercolorbox}[sep=12pt,center]{section title}
    \usebeamerfont{section title}\insertsection\par
  \end{beamercolorbox}
}
\setbeamertemplate{subsection page}{
  \centering
  \begin{beamercolorbox}[sep=8pt,center]{subsection title}
    \usebeamerfont{subsection title}\insertsubsection\par
  \end{beamercolorbox}
}
% Prevent slide breaks in the middle of a paragraph
\widowpenalties 1 10000
\raggedbottom
\AtBeginPart{
  \frame{\partpage}
}
\AtBeginSection{
  \ifbibliography
  \else
    \frame{\sectionpage}
  \fi
}
\AtBeginSubsection{
  \frame{\subsectionpage}
}
\usepackage{iftex}
\ifPDFTeX
  \usepackage[T1]{fontenc}
  \usepackage[utf8]{inputenc}
  \usepackage{textcomp} % provide euro and other symbols
\else % if luatex or xetex
  \usepackage{unicode-math} % this also loads fontspec
  \defaultfontfeatures{Scale=MatchLowercase}
  \defaultfontfeatures[\rmfamily]{Ligatures=TeX,Scale=1}
\fi
\usepackage{lmodern}
\ifPDFTeX\else
  % xetex/luatex font selection
\fi
% Use upquote if available, for straight quotes in verbatim environments
\IfFileExists{upquote.sty}{\usepackage{upquote}}{}
\IfFileExists{microtype.sty}{% use microtype if available
  \usepackage[]{microtype}
  \UseMicrotypeSet[protrusion]{basicmath} % disable protrusion for tt fonts
}{}
\makeatletter
\@ifundefined{KOMAClassName}{% if non-KOMA class
  \IfFileExists{parskip.sty}{%
    \usepackage{parskip}
  }{% else
    \setlength{\parindent}{0pt}
    \setlength{\parskip}{6pt plus 2pt minus 1pt}}
}{% if KOMA class
  \KOMAoptions{parskip=half}}
\makeatother
\usepackage{longtable,booktabs,array}
\usepackage{caption}
\captionsetup[table]{skip=6pt}
\newcounter{none} % for unnumbered tables
\usepackage{calc} % for calculating minipage widths
\usepackage{caption}
% Make caption package work with longtable
\makeatletter
\def\fnum@table{\tablename~\thetable}
\makeatother
\setlength{\emergencystretch}{3em} % prevent overfull lines
\providecommand{\tightlist}{%
  \setlength{\itemsep}{0pt}\setlength{\parskip}{0pt}}
% Auto-generated by infrastructure/rendering/_pdf_latex_helpers.py::extract_math_font_preamble.
% Minimal header passed to Pandoc via -H so Beamer slides pick up the same math font
% as the combined-PDF path. Adjust the manuscript's preamble.md to change the font globally.
\usepackage{unicode-math}
\setmathfont{latinmodern-math.otf}

% Slide layout defaults for warning-clean scientific decks.
\usepackage{etoolbox}
\IfFileExists{xurl.sty}{\usepackage{xurl}}{}
\IfFileExists{seqsplit.sty}{\usepackage{seqsplit}}{\newcommand{\seqsplit}[1]{#1}}
\protected\def\breakseq#1{\seqsplit{#1}}
\protected\def\breaktt#1{\begingroup\ttfamily\seqsplit{#1}\endgroup}
\setlength{\emergencystretch}{6em}
\tolerance=5000
\hbadness=10000
\hfuzz=1pt
\setlength{\tabcolsep}{2pt}
\AtBeginEnvironment{longtable}{\tiny\renewcommand{\arraystretch}{0.86}\setlength{\tabcolsep}{1pt}}
\AtBeginEnvironment{tabular}{\tiny\renewcommand{\arraystretch}{0.86}\setlength{\tabcolsep}{1pt}}
\AtBeginEnvironment{equation}{\tiny}
\AtBeginEnvironment{equation*}{\tiny}
\AtBeginEnvironment{align}{\tiny}
\AtBeginEnvironment{align*}{\tiny}
\AtBeginEnvironment{itemize}{\footnotesize}
\AtBeginEnvironment{enumerate}{\footnotesize}
\AtBeginEnvironment{description}{\footnotesize}
\setbeamerfont{caption}{size=\tiny}
\setbeamerfont{caption name}{size=\tiny}
\setbeamerfont{normal text}{size=\small}
\setbeamerfont{frametitle}{size=\small}
\setbeamerfont{section title}{size=\footnotesize}
\setbeamerfont{subsection title}{size=\footnotesize}
\setbeamertemplate{section page}{%
  \centering
  \begin{beamercolorbox}[sep=12pt,center,wd=\paperwidth]{section title}
    \parbox{0.86\paperwidth}{\centering\usebeamerfont{section title}\insertsection\par}
  \end{beamercolorbox}
}
\setbeamertemplate{subsection page}{%
  \centering
  \begin{beamercolorbox}[sep=8pt,center,wd=\paperwidth]{subsection title}
    \parbox{0.86\paperwidth}{\centering\usebeamerfont{subsection title}\insertsubsection\par}
  \end{beamercolorbox}
}
\setlength{\abovecaptionskip}{2pt}
\setlength{\belowcaptionskip}{0pt}

% Natbib and cross-reference fallbacks — slides don't load natbib
% or cleveref, but combined-PDF manuscript prose may emit these
% commands. The fallback renders citations as a bracketed key list
% and cross-references as detokenized labels so slides don't fail on
% undefined control sequences or raw underscores. \providecommand is
% a no-op if packages load later.
\providecommand{\citep}[1]{[#1]}
\providecommand{\citet}[1]{#1}
\providecommand{\citealp}[1]{#1}
\providecommand{\citeauthor}[1]{#1}
\providecommand{\citeyear}[1]{#1}
\providecommand{\cref}[1]{\texttt{\detokenize{#1}}}
\providecommand{\Cref}[1]{\texttt{\detokenize{#1}}}

\newtheorem{proposition}{Proposition}
\newtheorem{hypothesis}{Hypothesis}
\usepackage{bookmark}
\IfFileExists{xurl.sty}{\usepackage{xurl}}{} % add URL line breaks if available
\urlstyle{same}
% fallback for those not using the hyperref driver hyperxmp:
\makeatletter
\@ifundefined{xmpquote}{\newcommand{\xmpquote}[1]{#1}}{}
\makeatother
\hypersetup{
  hidelinks,
  pdfcreator={LaTeX via pandoc}}

\author{\texorpdfstring{}{}}
\date{}

\begin{document}

\section{Limitations and Planned Evaluation: Delivery Dependencies,
Kernel Scope, and Learner
Outcomes}\label{sec:limitations-delivery-dependencies-learner-outcomes}

\subsection{Delivery dependencies, fixed-kernel scope, and
learner-outcome
boundaries}\label{delivery-dependencies-fixed-kernel-scope-and-learner-outcome-boundaries}

\begin{frame}[fragile,allowframebreaks]{Delivery substrate and kernel
scope}
\protect\phantomsection\label{delivery-substrate-and-kernel-scope}
The local learner dashboard can be verified in artifact-only mode
without Docker; live progress, ranking, and quiz completion still
require a SkillTree service. When that service is available, the
dashboard mounts the official learner display and separates local lesson
material from official subject and quiz handoffs, but it does not
provide a second source of learner truth. The REST contract is verified
against a \breaktt{pytest-httpserver} mirror of SkillTree's endpoints,
and the live seed is documented and scripted.

\end{frame}

\begin{frame}[fragile,allowframebreaks]{Delivery substrate and kernel
scope}
The implemented kernels cover discrete pedagogical cases, reference
variational examples, and bounded predictive-coding or planning
demonstrations; they do not provide a continuous, scalable, robust, or
deep predictive-coding inference engine. Reactive message passing,
RxInfer, coupled free energy, active predictive coding,
divide-and-conquer predictive coding, generalized filtering, generalized
free energy, hierarchical multi-step planning, and richer
reinforcement-learning comparisons remain cited roadmap context rather
than hidden implementation claims (Bagaev and Vries 2023; Bagaev et al.
2023; Nelson et al. 2025; Rao et al. 2024; Sennesh et al. 2024; Friston
et al. 2008; Friston et al. 2010; Parr and Friston 2019). The RL
comparison matters because probabilistic-inference views of control
overlap with, but are not identical to, active-inference policy
selection (Levine 2018). Psychology-and-psychiatry Active Inference
\end{frame}

\begin{frame}[fragile,allowframebreaks]{Delivery substrate and kernel
scope}
reviews likewise support cautious application language rather than
theory-level validation claims (Badcock and Davey 2024). Finally, the
platform's contribution is the verifiability and provenance of
curriculum content, not a measured claim about learning outcomes: it
does not evaluate teaching effectiveness, retention, or transfer, and no
efficacy study is reported here. Microlearning studies,
epistemic-curiosity evidence, learning-analytics reviews, and
learning-analytics intervention meta-analyses motivate the proposed
study design but do not substitute for data from this curriculum, and
learning analytics also introduces privacy and data-protection
obligations before progress exports can be treated as evidence (Monib et
al. 2025; Rof et al. 2024; Emin et al. 2026; Liu et al. 2025;
Hernández-Campos et al. 2025; Liu and Khalil 2023). Whether this
provenance discipline actually improves how people learn the framework
\end{frame}

\begin{frame}[fragile,allowframebreaks]{Delivery substrate and kernel
scope}
remains future work, outside the scope of what the code, citations, and
validators establish. The preregistered pilot design in
\breaktt{experiment\_plan.yaml} is therefore a setup artifact: it defines
consent, privacy, SkillTree progress-export, pre/post,
delayed-retention, and analysis boundaries, and follows design-based
research framing for studying a curriculum artifact in context, but it
is not learner-outcome evidence until real data are collected and added
to the claim-support record (Design-Based Research Collective 2003;
Barab and Squire 2004).
\end{frame}

\end{document}
